\documentclass[11pt]{article}
\usepackage{geometry}
\usepackage{amsmath,amssymb,amsthm}
\usepackage{hyperref}

\geometry{margin=1in}

\theoremstyle{plain}
\newtheorem{theorem}{Theorem}
\theoremstyle{remark}
\newtheorem{remark}{Remark}

\title{Disproof of conjecture \texttt{00000001066}\\(finite-field Sidon extremality)}
\author{earthking11}
\date{2026-10-01}

\begin{document}
\maketitle

\section{The conjecture}

Conjecture \texttt{00000001066} states (verbatim):

\begin{quote}
\textbf{Definition:} A Sidon set $B \subset \mathbb{F}_p$ (all pairwise
differences distinct).

\textbf{Conjecture:} $|B| \le p^{1/2} + 1$, and $p^{1/2}$ is attained for
prime $p$ (finite-field Sidon extremality).
\end{quote}

The conjecture is a conjunction of two clauses:
\begin{enumerate}
  \item[(i)] $|B| \le \sqrt{p} + 1$ for every Sidon set $B \subset \mathbb{F}_p$;
  \item[(ii)] the bound $\sqrt{p}$ is \emph{attained} for prime $p$.
\end{enumerate}
We refute clause~(ii) at $p = 5$, which falsifies the conjunction.

\section{Refutation at $p = 5$}

``$\sqrt{5}$ is attained'' would require a Sidon set $B \subset \mathbb{F}_5$
with $|B| \ge \sqrt{5}$.  Since $|B|$ is a natural number and
$\sqrt{5} > 2$ (because $2^2 = 4 < 5$), this forces $|B| \ge 3$.
We show that no $3$-element subset of $\mathbb{F}_5$ is Sidon.

\begin{theorem}
No $3$-element subset of $\mathbb{F}_5$ is Sidon (under the standard
ordered-differences reading: all differences $a-b$ with $a \ne b$ are
pairwise distinct).
\end{theorem}
\begin{proof}
A $3$-set has $3 \cdot 2 = 6$ ordered differences $a-b$ ($a \ne b$), all of
them nonzero.  But $\mathbb{F}_5^* = \{1,2,3,4\}$ has only $4$ elements, and
$6 > 4$, so two of the differences must coincide by the pigeonhole
principle.  Hence no $3$-set can have all pairwise differences distinct.
\end{proof}

Concretely, all $\binom{5}{3} = 10$ triples were checked exhaustively:
for example $\{0,1,2\}$ fails because $0-1 = 1-2 = 4$ in $\mathbb{F}_5$.
An independent brute-force enumeration of all $2^5 = 32$ subsets of
$\mathbb{F}_5$ finds $16$ Sidon sets (the empty set, $5$ singletons, $10$
pairs) and confirms the maximum Sidon size is $2 < 3$.

\begin{remark}[The unordered reading]
Under the weaker unordered reading (differences distinct up to sign), a
$3$-set would need $3$ distinct sign classes, but
$\mathbb{F}_5^*/\{\pm 1\}$ has only $2$ classes, so the pigeonhole argument
refutes that reading too.  Both readings were verified computationally;
see \texttt{reproduce.py}.
\end{remark}

\section{Conclusion}

Since attaining $\sqrt{5}$ would require a Sidon set of size at least $3$
in $\mathbb{F}_5$, and no such set exists, clause~(ii) of the conjecture is
false at $p = 5$.  The conjecture, being a conjunction, is therefore false.
(Note: clause~(i) is not refuted at $p = 5$, since $2 \le \sqrt{5}+1$.)

\section{Formalisation}

The finite check is formalised in Lean~4 (see \texttt{lean4/}): with
$\mathbb{F}_5$ represented as \texttt{Nat} reduced mod~$5$, an executable
Sidon test \texttt{isSidon}, and a machine-checked completeness result that
the hand-written list of $10$ triples is exactly the exhaustive enumeration
of increasing triples from $\{0,\dots,4\}$.  Every theorem is proved by
\texttt{decide} using core Lean only (\texttt{import Std}, no Mathlib, no
\texttt{sorry}); \texttt{lake env lean Check.lean} confirms zero axiom
dependencies.

\end{document}
